% ===== Chapter 8 opening: what counts as a mental disorder (~8 min) =====
\partopener{Գլուխ 8}{Երբ ուղեղը փչանում է}{Հոգեկան առողջության խնդիրներ, և ինչպես են դրանք առաջանում}{thinking}{lavender}

% Neurological vs mental: hardware vs software
\begin{frame}{Երկու տեսակի «փչացում»}
  \begin{columns}[T]
    \begin{column}{0.47\textwidth}
      \begin{card}[colback=peach!35!cardwhite]
        \emoji[1.3em]{plug}\ {\blackfont Նյարդաբանական}\par\smallskip
        Ֆիզիկական վնաս ուղեղում.
        \begin{itemize}\setlength{\itemsep}{0.2mm}
          \item հիպոկամպ $\rightarrow$ ամնեզիա
          \item սև նյութ $\rightarrow$ Պարկինսոն
          \item նոպաներ, միգրեն
        \end{itemize}
      \end{card}
    \end{column}
    \begin{column}{0.47\textwidth}
      \onslide<2->{%
      \begin{card}[colback=lavender!40!cardwhite]
        \emoji[1.3em]{floppy}\ {\blackfont Հոգեկան}\par\smallskip
        Մտածողության, վարքի կամ զգացմունքների շեղում։ Ուղեղը ֆիզիկապես \textbf{նորմալ} է, բայց անում է ոչ օգտակար բաներ
      \end{card}}
    \end{column}
  \end{columns}
  \vspace{0.5mm}
  \onslide<3->{%
  \begin{bookquote}
    Նյարդաբանականը «երկաթի» խնդիր է, հոգեկանը՝ «ծրագրի»։ Թեև սահմանն այնքան էլ հստակ չէ։
  \end{bookquote}}
\end{frame}

% No two brains are alike: who is "normal"?
\begin{frame}{Ո՞ր ուղեղն է «նորմալ»}
  \begin{columns}[c]
    \begin{column}{0.52\textwidth}
      \begin{itemize}
        \item Միլիարդավոր նեյրոններ, տրիլիոնավոր կապեր. երկու միանման ուղեղ չկա
        \item Ամեն մեկս ունի իր տարօրինակությունները
        \item<2-> \textbf{Սինեսթեզիա} (synaesthesia). շատերը չգիտեն, որ ունեն, մինչև չասեն՝ «սիրում եմ մանուշակագույնի հոտը»
      \end{itemize}
    \end{column}
    \begin{column}{0.42\textwidth}
      \centering
      \onslide<2->{\bigemoji[1.1cm]{palette}\ \raisebox{0.3cm}{\Large +}\ \bigemoji[1.1cm]{nose}\par\medskip}
      \onslide<3->{%
      \begin{card}[colback=butter!45!cream]
        {\blackfont Հոգեկան խանգարում}՝ մտքի կամ վարքի նախշ, որը \textbf{տառապանք} է պատճառում կամ խանգարում է գործել «նորմալ» հասարակության մեջ
      \end{card}}
    \end{column}
  \end{columns}
\end{frame}

% "Normal" moves with time
\begin{frame}{«Նորմալը» փոխվում է}
  \begin{columns}[c]
    \begin{column}{0.4\textwidth}
      \begin{card}[colback=brainpink!25!cardwhite]
        \centering
        {\blackfont\Large 1973}\par\smallskip
        \raggedright Միայն այդ տարի Ամերիկյան հոգեբուժական ասոցիացիան համասեռականությունը հանեց խանգարումների ցանկից
      \end{card}
    \end{column}
    \begin{column}{0.54\textwidth}
      \onslide<2->{%
      Ցանկը անընդհատ վերանայվում է.
      \begin{itemize}\setlength{\itemsep}{0.4mm}
        \item նոր գիտելիք, նոր բուժումներ
        \item գիտական «դպրոցների» փոփոխություն
        \item դեղագործական ընկերություններ, որոնք սիրում են նոր հիվանդություններ՝ նոր դեղեր վաճառելու համար
      \end{itemize}}
    \end{column}
  \end{columns}
  \vspace{3mm}
  \onslide<3->{%
  \begin{card}[colback=lavender!35!cardwhite]
    \centering «Խանգարման» ու «նորմայի» սահմանը շատ \textbf{մշուշոտ} է ու հաճախ կախված է հասարակության կանոններից
  \end{card}}
\end{frame}

% How common: guess first, then the book's number
\begin{frame}{Որքա՞ն տարածված է}
  \begin{columns}[c]
    \begin{column}{0.4\textwidth}
      \begin{chatbox}
        Գուշակիր. 100 մարդուց քանի՞սն են կյանքում ունենում հոգեկան խանգարման որևէ դրսևորում։
      \end{chatbox}
      \vspace{1mm}
      \onslide<3->{%
      \begin{card}[colback=mint!30!cardwhite]
        Ավելի նոր տվյալներ. կյանքում՝ \textbf{29\%}, վերջին տարում՝ \textbf{18\%}
      \end{card}}
    \end{column}
    \begin{column}{0.56\textwidth}
      \centering
      \onslide<2->{%
      {\blackfont\Large Գրեթե \textcolor{pinkdark}{4-ից 1-ը}}\par\medskip
      \begin{tikzpicture}
        \foreach \i in {0,...,19}{
          \pgfmathsetmacro{\px}{mod(\i,10)*0.66}
          \pgfmathsetmacro{\py}{-floor(\i/10)*1.15}
          \pgfmathtruncatemacro{\hl}{mod(\i*7+3,4)==0 ? 1 : 0}
          \ifnum\hl=1 \def\pc{brainpink}\else\def\pc{lavender!45!cardwhite}\fi
          \filldraw[fill=\pc, draw=plum, line width=0.7pt] (\px,\py+0.52) circle (1.6mm);
          \filldraw[fill=\pc, draw=plum, line width=0.7pt, rounded corners=1.6mm]
            (\px-0.24,\py-0.08) rectangle (\px+0.24,\py+0.32);
        }
      \end{tikzpicture}\par\medskip
      {\small Բայց երբ քեզ հետ չէ, հեշտ է \textbf{անտեսել}}}
    \end{column}
  \end{columns}
  \source{Steel et al., Int. J. Epidemiol., 2014}
\end{frame}

% What to call it
\begin{frame}{Ինչպե՞ս անվանել}
  \begin{columns}[T]
    \begin{column}{0.31\textwidth}
      \begin{card}[colback=peach!35!cardwhite]
        {\blackfont «Հիվանդություն»}\par\smallskip
        Հնչում է գրիպի պես. բուժվում է ու անցնում։ Բայց հաճախ «նորոգելու» բան չկա
      \end{card}
    \end{column}
    \begin{column}{0.31\textwidth}
      \onslide<2->{%
      \begin{card}[colback=lavender!40!cardwhite]
        {\blackfont «Խանգարում»}\par\smallskip
        Հնչում է որպես «վնասված»։ Գուցե դա մտածելու \textbf{այլ ձև} է
      \end{card}}
    \end{column}
    \begin{column}{0.31\textwidth}
      \onslide<3->{%
      \begin{card}[colback=mint!35!cardwhite]
        {\blackfont Չեզոք բառեր}\par\smallskip
        Շատ հոգեբաններ կարծում են, որ «հիվանդ» պիտակն ինքնին վնասում է
      \end{card}}
    \end{column}
  \end{columns}
  \vspace{3mm}
  \onslide<4->{%
  \begin{card}[colback=butter!45!cream]
    \emoji{stethoscope}\ Բըրնեթը մնում է \textbf{բժշկական} տեսանկյունին. դա իր մասնագիտությունն է ու ամենածանոթ լեզուն
  \end{card}}
\end{frame}

% Why the chapter exists, and today's tone
\begin{frame}{Ինչի՞ համար է այս գլուխը}
  \begin{columns}[c]
    \begin{column}{0.62\textwidth}
      \begin{bookquote}
        Բացատրել, թե ինչպես է ուղեղը մեզ «մատնում»։ Ե՛վ նրանց, ում հետ դա պատահում է, և՛ նրանց շրջապատին, որ հաճախ չի հասկանում՝ ինչ է կատարվում։
      \end{bookquote}
      \vspace{2mm}
      \onslide<2->{%
      \begin{card}[colback=mint!30!cardwhite]
        Այսօր էլ կկատակենք, բայց \textbf{ուղեղի} ու Բըրնեթի մասին։ Ոչ երբեք նրանց մասին, ում հետ սա պատահում է։
      \end{card}}
    \end{column}
    \begin{column}{0.32\textwidth}
      \centering\mascot[2.8cm]{thinking}
    \end{column}
  \end{columns}
\end{frame}
